\documentclass[11pt,a4paper]{article}
\usepackage[T1]{fontenc}
\usepackage[utf8]{inputenc}
\usepackage[english]{babel}
\usepackage{lmodern}
\usepackage{microtype}
\usepackage{geometry}
\geometry{margin=2.5cm}
\usepackage{hyperref}
\hypersetup{
    colorlinks=true,
    linkcolor=blue!70!black,
    urlcolor=blue!70!black,
    citecolor=blue!70!black
}
\usepackage{booktabs}
\usepackage{longtable}
\usepackage{enumitem}
\usepackage{xcolor}
\usepackage{titlesec}
\usepackage{fancyhdr}
\usepackage{graphicx}
\usepackage{listings}

\definecolor{codegray}{rgb}{0.95,0.95,0.95}
\definecolor{doctum}{rgb}{0.10,0.30,0.60}

\lstset{
  backgroundcolor=\color{codegray},
  basicstyle=\small\ttfamily,
  breaklines=true,
  frame=single,
  rulecolor=\color{gray!40},
  xleftmargin=1em,
  framexleftmargin=0.5em
}

\pagestyle{fancy}
\fancyhf{}
\fancyhead[L]{\textcolor{doctum}{\small\textbf{Doctum Consilium}}}
\fancyhead[R]{\small\nouppercase{\leftmark}}
\fancyfoot[C]{\thepage}

\titleformat{\section}{\large\bfseries\color{doctum}}{}{0em}{}[\titlerule]
\titleformat{\subsection}{\normalsize\bfseries\color{doctum!80}}{}{0em}{}


\begin{document}

\title{
  \textcolor{doctum}{\Large\textbf{Guardrails Kit}} \\[0.5em]
  \large Multi-Repo AI Assistant Policy and Engineering Standards Framework \\[0.3em]
  \normalsize Technical Foundations and Architecture
}
\author{Doctum Consilium \\ \small\textit{Generated 2026-05-08}}
\date{}
\maketitle
\thispagestyle{fancy}

\begin{abstract}
Canonical guardrails, templates, and agents for 13 downstream repositories. Ensures consistent code documentation, secret management, and deployment standards across GitHub Copilot, Claude Code, and Gemini CLI.
This document covers the theoretical foundations, architectural decisions,
and key technical concepts required to understand, contribute to, and
maintain this repository.
\end{abstract}

\tableofcontents
\newpage

\section{Overview}

Guardrails Kit is a meta-repository that defines and enforces engineering standards
across an entire GitHub organisation. It provides:
\begin{itemize}
  \item Policy documents in \texttt{.guardrails/rules/}
  \item AI assistant configuration templates (Copilot, Claude, Gemini)
  \item A Bash/PowerShell validator script for CI enforcement
  \item GitHub Actions workflow for automated gate checks
  \item Installation scripts for downstream repos
\end{itemize}
The multi-AI design ensures that whether a contributor uses GitHub Copilot,
Claude Code, or Gemini CLI, the same guardrails are applied.


\section{Architecture}

\begin{verbatim}
guardrails-kit (canonical source)
├── templates/
│   ├── copilot-instructions.template.md  → .github/copilot-instructions.md
│   ├── CLAUDE.template.md                → CLAUDE.md
│   └── GEMINI.template.md                → GEMINI.md
├── .guardrails/
│   ├── rules/*.md         (policy documents)
│   └── bin/validate_guardrails.sh
└── scripts/
    ├── install_into_repo.sh   (Ubuntu/macOS)
    └── install_into_repo.ps1  (Windows)

Downstream repos (×13)
├── .github/copilot-instructions.md
├── CLAUDE.md
├── GEMINI.md
└── .github/agents/ecr-secrets-agent.agent.md
\end{verbatim}


\section{Key Technical Concepts}
\begin{itemize}[leftmargin=1.5em]
  \item Infrastructure-as-Policy: treating engineering standards as versioned, testable artefacts rather than wiki pages.
  \item AI assistant instruction files: \texttt{.github/copilot-instructions.md} is loaded by GitHub Copilot; \texttt{CLAUDE.md} by Claude Code; \texttt{GEMINI.md} by Gemini CLI at session start.
  \item CI gate pattern: \texttt{validate\_guardrails.sh} checks required docs, doc-update rules, and optional lint/test/build commands as a pre-merge gate.
  \item Template propagation: \texttt{install\_into\_repo.sh} substitutes \texttt{\{\{PROJECT\_NAME\}\}} placeholder and copies files; idempotent with \texttt{--force} flag.
  \item Platform compatibility: all scripts provide Ubuntu (Bash), Windows (PowerShell), and macOS (zsh/bash) variants.
\end{itemize}

\section{Data Flow}

guardrails-kit template update →
\texttt{scripts/install\_into\_repo.sh /path/to/repo} →
file copy + placeholder substitution →
\texttt{git add + commit} in target repo →
GitHub Actions \texttt{Guardrails} check passes.


\section{Technology Stack}

\begin{tabular}{ll}
\toprule
\textbf{Component} & \textbf{Technology} \\
\midrule
Shell scripting & Bash 5, PowerShell 7 \\
CI & GitHub Actions \\
AI assistants & GitHub Copilot, Claude Code, Gemini CLI \\
Markdown & CommonMark (instruction files) \\
LaTeX & pdflatex (documentation) \\
\bottomrule
\end{tabular}


\section{Design Decisions}

\begin{itemize}
  \item Single source of truth: all 13 repos pull standards from one place;
    a policy change propagates to all repos in one \texttt{install} pass.
  \item Three AI file targets: avoids vendor lock-in; contributors can use
    whichever AI assistant they prefer without losing guardrail enforcement.
  \item Validator checks staged files: integrates cleanly with pre-commit hooks
    and GitHub Actions \texttt{on: pull\_request} trigger.
\end{itemize}


\section{Repository Structure}
\begin{lstlisting}[language=bash,caption={Top-level directory layout}]
guardrails-kit/
  .github/
    copilot-instructions.md   # GitHub Copilot guardrails
    agents/
      ecr-secrets-agent.agent.md  # ECR secret rotation runbook
  CLAUDE.md                   # Claude Code guardrails
  GEMINI.md                   # Gemini CLI guardrails
  INFRASTRUCTURE.md           # Operational runbook
  ROADMAP.md                  # Execution log and roadmap
  README.md                   # Project overview
\end{lstlisting}

\section{Contributing}
\begin{enumerate}
  \item Read \texttt{README.md}, \texttt{INFRASTRUCTURE.md}, and \texttt{ROADMAP.md} before any task.
  \item Follow the Code Documentation Standard: every public function must have
    a docstring (Google-style for Python, JSDoc for TypeScript, XML for C\#).
  \item Run the guardrails validator before committing:
    \begin{lstlisting}[language=bash]
git add -A && ./.guardrails/bin/validate_guardrails.sh
    \end{lstlisting}
  \item For deployment or destructive operations, always use a named tmux session:
    \begin{lstlisting}[language=bash]
tmux new-session -A -s deploy-guardrails-kit
    \end{lstlisting}
\end{enumerate}

\end{document}
